% !PRINTED 530
\weberpage{530}
This $\varDelta$ is a positive or negative, but never vanishing, integral rational number completely determined by $\Omega$, and there is always a basis of $\Omega$ consisting of integers $\omega_1, \omega_2, \ldots, \omega_n$ whose discriminant is equal to $\varDelta$.

Such a basis we shall briefly call a \gesperrt{minimal basis of} $\Omega$.

If we understand by $k_1, k_2, \ldots, k_n$ any integral rational numbers whatever, then every number of the shape
\weq{1}{\omega = k_1\omega_1 + k_2\omega_2 + \cdots + k_n\omega_n}
is an integral algebraic number, and we now prove the fundamental theorem:

1. \gesperrt{If $\omega_1, \omega_2, \ldots, \omega_n$ is a minimal basis, then all the integers of the field $\Omega$ are contained in the form (1).}

Since $\omega_1, \omega_2, \ldots, \omega_n$ is a basis, every number in $\Omega$ can first of all be represented in the form (1) if rational fractions are admitted for $k_1, k_2, \ldots, k_n$. Let us assume, then, that an integer $\omega$ is representable in the form
\weq{2}{\omega = \frac{k_1\omega_1 + k_2\omega_2 + \cdots + k_n\omega_n}{k}}
where $k_1, k_2, \ldots, k_n, k$ are integral rational numbers such that not all of $k_1, k_2, \ldots, k_n$ have a common divisor with $k$. If $p$ is any natural prime dividing $k$, and $k = pk'$, then at least one of the coefficients $k_1, k_2, \ldots, k_n$ must be indivisible by $p$. Suppose, say, that $k_1$ is not divisible by $p$; then the integral rational number $l$ can be determined so that $l k_1 \equiv 1 \; (\mathrm{mod}\ p)$, or $(l k_1 - 1)$ becomes divisible by $p$. It then follows from (2):
\weq{3}{l k' \omega - \frac{l k_1 - 1}{p}\omega_1 = \frac{\omega_1 + l k_2\omega_2 + \cdots + l k_n\omega_n}{p} = \omega_1',}
and $\omega_1'$ is likewise an integral algebraic number. If we set
\weq{4}{\omega_2' = \omega_2, \ldots, \omega_n' = \omega_n,}
then the numbers $\omega_1', \omega_2', \ldots, \omega_n'$ form an \gesperrt{integral basis} of $\Omega$, because the numbers $\omega_1, \omega_2, \ldots, \omega_n$, and consequently all the numbers $\omega$, can be expressed linearly in terms of $\omega_1', \omega_2', \ldots, \omega_n'$, and formula (11) of the preceding section yields
\weq{5}{\varDelta(\omega_1', \omega_2', \ldots, \omega_n') = \frac{1}{p^2}\varDelta(\omega_1, \omega_2, \ldots, \omega_n).}
